% GENERATED FILE. Edit canonical lessons, then run npm run generate:books.
% canonical-source-hash: fnv1a64:a131cd90a3566a5c
% canonical-lessons: SA-C340-nimittam, SA-C340-prabhavah, SA-C340-abhavah, SA-C340-utpadyate, SA-C340-janayati

\chapter{Why It Happened}
\label{ch:sa-why-it-happened}
\hlchaptermodality{340}

\begin{chapteropening}
\textbf{By the end of this chapter:} \emph{I can say what caused something and what it led to.}

Complete the last lesson of chapter 340: I can say what caused something and what it led to.
\end{chapteropening}

\section[\texorpdfstring{\sk{निमित्तम्} (nimittam)}{निमित्तम् (nimittam)}]{\sk{निमित्तम्} (nimittam) — an occasion; a cause}

\label{lesson:SA-C340-nimittam}



\begin{quote}
Before the new one: say the Sanskrit for radio, then the Sanskrit for the internet.
\end{quote}

\subsection*{You'll want to know: \sk{निमित्तम्}}
\textbf{\sk{निमित्तम्}} — \emph{nimittam} — \textquotedblleft{}an occasion; a cause\textquotedblright{}.

\begin{quote}
\sk{विवाहः} \sk{मम} \sk{आगमनस्य} \sk{निमित्तम्} \sk{अस्ति।}
\end{quote}

\begin{quote}
\emph{vivāhaḥ mama āgamanasya nimittam asti.}
\end{quote}

\begin{quote}
A wedding is the occasion of my visit.
\end{quote}

A neuter noun.

\begin{cousinweb}
\textbf{\sk{निमित्तम्}} (\emph{nimittam}) MW gives \emph{nimitta} the senses \textquotedblleft{}a mark\textquotedblright{}, \textquotedblleft{}a sign\textquotedblright{} and \textquotedblleft{}a cause, a motive, a ground\textquotedblright{}, and says it is possibly connected with \emph{ni} and the root √mā, \textquotedblleft{}to measure\textquotedblright{}; so its origin is uncertain (MW).
\end{cousinweb}

\begin{grammarlens}[title={What you've built}]
One more word for giving reasons and opinions.
\end{grammarlens}

\subsection*{Your turn}
\begin{itemize}
\raggedright
  \item \emph{Say it:} \emph{nimittam}
  \item \emph{Say it:} \emph{nimittam}, once more
  \item \emph{Say it:} \emph{vivāhaḥ mama āgamanasya nimittam asti.} — A wedding is the occasion of my visit.
  \item \emph{Recall:} say the Sanskrit for radio, then the Sanskrit for the internet, then say \emph{nimittam} again
\end{itemize}

\begin{tcolorbox}[breakable,colback=teal!4,colframe=teal!35!black,title={Before you move on}]
What does \sk{निमित्तम्} mean? (\textquotedblleft{}An occasion; a cause\textquotedblright{}.) Say it once more.
\end{tcolorbox}
\section[\texorpdfstring{\sk{प्रभावः} (prabhāvaḥ)}{प्रभावः (prabhāvaḥ)}]{\sk{प्रभावः} (prabhāvaḥ) — an influence; power}

\label{lesson:SA-C340-prabhavah}



\begin{quote}
Before the new one: say the Sanskrit for the internet, then the Sanskrit for an occasion.
\end{quote}

\subsection*{You'll want to know: \sk{प्रभावः}}
\textbf{\sk{प्रभावः}} — \emph{prabhāvaḥ} — \textquotedblleft{}an influence; power\textquotedblright{}.

\begin{quote}
\sk{अध्यापकस्य} \sk{प्रभावः} \sk{छात्रेषु} \sk{महान्} \sk{अस्ति।}
\end{quote}

\begin{quote}
\emph{adhyāpakasya prabhāvaḥ chātreṣu mahān asti.}
\end{quote}

\begin{quote}
The teacher's influence on the students is great.
\end{quote}

A masculine noun.

\begin{cousinweb}
\textbf{\sk{प्रभावः}} (\emph{prabhāvaḥ}) is built on \emph{pra}, \textquotedblleft{}forth\textquotedblright{}, and the root √bhū, \textquotedblleft{}to be, to become\textquotedblright{}; MW gives \emph{pra-bhāva} the senses \textquotedblleft{}might, power, efficacy\textquotedblright{} (MW).
\end{cousinweb}

\begin{grammarlens}[title={What you've built}]
One more word for giving reasons and opinions.
\end{grammarlens}

\subsection*{Your turn}
\begin{itemize}
\raggedright
  \item \emph{Say it:} \emph{prabhāvaḥ}
  \item \emph{Say it:} \emph{prabhāvaḥ}, once more
  \item \emph{Say it:} \emph{adhyāpakasya prabhāvaḥ chātreṣu mahān asti.} — The teacher's influence on the students is great.
  \item \emph{Recall:} say the Sanskrit for the internet, then the Sanskrit for an occasion, then say \emph{prabhāvaḥ} again
\end{itemize}

\begin{tcolorbox}[breakable,colback=teal!4,colframe=teal!35!black,title={Before you move on}]
What does \sk{प्रभावः} mean? (\textquotedblleft{}An influence; power\textquotedblright{}.) Say it once more.
\end{tcolorbox}
\section[\texorpdfstring{\sk{अभावः} (abhāvaḥ)}{अभावः (abhāvaḥ)}]{\sk{अभावः} (abhāvaḥ) — a lack; an absence}

\label{lesson:SA-C340-abhavah}



\begin{quote}
Before the new one: say the Sanskrit for an occasion, then the Sanskrit for an influence.
\end{quote}

\subsection*{You'll want to know: \sk{अभावः}}
\textbf{\sk{अभावः}} — \emph{abhāvaḥ} — \textquotedblleft{}a lack; an absence\textquotedblright{}.

\begin{quote}
\sk{जलस्य} \sk{अभावः} \sk{अस्ति}, \sk{अतः} \sk{वृक्षाः} \sk{शुष्यन्ति।}
\end{quote}

\begin{quote}
\emph{jalasya abhāvaḥ asti, ataḥ vṛkṣāḥ śuṣyanti.}
\end{quote}

\begin{quote}
There is a lack of water, so the trees are drying up.
\end{quote}

A masculine noun.

\begin{cousinweb}
\textbf{\sk{अभावः}} (\emph{abhāvaḥ}) is built from \emph{a-}, \textquotedblleft{}not\textquotedblright{}, and \emph{bhāva}, \textquotedblleft{}being\textquotedblright{}, from the root √bhū; MW gives \emph{a-bhāva} the senses \textquotedblleft{}non-existence, absence\textquotedblright{} (MW).
\end{cousinweb}

\begin{grammarlens}[title={What you've built}]
One more word for giving reasons and opinions.
\end{grammarlens}

\subsection*{Your turn}
\begin{itemize}
\raggedright
  \item \emph{Say it:} \emph{abhāvaḥ}
  \item \emph{Say it:} \emph{abhāvaḥ}, once more
  \item \emph{Say it:} \emph{jalasya abhāvaḥ asti, ataḥ vṛkṣāḥ śuṣyanti.} — There is a lack of water, so the trees are drying up.
  \item \emph{Recall:} say the Sanskrit for an occasion, then the Sanskrit for an influence, then say \emph{abhāvaḥ} again
\end{itemize}

\begin{tcolorbox}[breakable,colback=teal!4,colframe=teal!35!black,title={Before you move on}]
What does \sk{अभावः} mean? (\textquotedblleft{}A lack; an absence\textquotedblright{}.) Say it once more.
\end{tcolorbox}
\section[\texorpdfstring{\sk{उत्पद्यते} (utpadyate)}{उत्पद्यते (utpadyate)}]{\sk{उत्पद्यते} (utpadyate) — comes about; arises}

\label{lesson:SA-C340-utpadyate}



\begin{quote}
Before the new one: say the Sanskrit for an influence, then the Sanskrit for a lack.
\end{quote}

\subsection*{You'll want to know: \sk{उत्पद्यते}}
\textbf{\sk{उत्पद्यते}} — \emph{utpadyate} — \textquotedblleft{}comes about; arises\textquotedblright{}.

\begin{quote}
\sk{क्रोधात्} \sk{विवादः} \sk{उत्पद्यते।}
\end{quote}

\begin{quote}
\emph{krodhāt vivādaḥ utpadyate.}
\end{quote}

\begin{quote}
Quarrels come about from anger.
\end{quote}

Past (it …): \sk{विवादः} \sk{क्रोधात्} \sk{उत्पन्नवान्}, \emph{vivādaḥ krodhāt utpannavān}.

\begin{cousinweb}
\textbf{\sk{उत्पद्यते}} (\emph{utpadyate}) is built on \emph{ut}, \textquotedblleft{}up, out\textquotedblright{}, and the root √pad, \textquotedblleft{}to go, to fall\textquotedblright{}; MW gives \emph{ut-√pad} the senses \textquotedblleft{}to arise, originate, be born or produced\textquotedblright{} (MW).
\end{cousinweb}

\begin{grammarlens}[title={What you've built}]
One more word for giving reasons and opinions.
\end{grammarlens}

\subsection*{Your turn}
\begin{itemize}
\raggedright
  \item \emph{Say it:} \emph{utpadyate}
  \item \emph{Say it:} \emph{utpadyate}, once more
  \item \emph{Say it:} \emph{krodhāt vivādaḥ utpadyate.} — Quarrels come about from anger.
  \item \emph{Recall:} say the Sanskrit for an influence, then the Sanskrit for a lack, then say \emph{utpadyate} again
\end{itemize}

\begin{tcolorbox}[breakable,colback=teal!4,colframe=teal!35!black,title={Before you move on}]
What does \sk{उत्पद्यते} mean? (\textquotedblleft{}Comes about; arises\textquotedblright{}.) Say it once more.
\end{tcolorbox}
\section[\texorpdfstring{\sk{जनयति} (janayati)}{जनयति (janayati)}]{\sk{जनयति} (janayati) — gives rise to; produces}

\label{lesson:SA-C340-janayati}



\begin{quote}
Before the new one: say the Sanskrit for a lack, then the Sanskrit for comes about.
\end{quote}

\subsection*{You'll want to know: \sk{जनयति}}
\textbf{\sk{जनयति}} — \emph{janayati} — \textquotedblleft{}gives rise to; produces\textquotedblright{}.

\begin{quote}
\sk{क्रोधः} \sk{दुःखं} \sk{जनयति।}
\end{quote}

\begin{quote}
\emph{krodhaḥ duḥkhaṃ janayati.}
\end{quote}

\begin{quote}
Anger gives rise to sorrow.
\end{quote}

Past (he …): \sk{सः} \sk{भयं} \sk{जनितवान्}, \emph{saḥ bhayaṃ janitavān}.

\begin{cousinweb}
\textbf{\sk{जनयति}} (\emph{janayati}) comes from the root √jan, \textquotedblleft{}to be born\textquotedblright{}, in its causative form, \textquotedblleft{}to cause to be born, to produce\textquotedblright{}; MW gives \emph{janita} \textquotedblleft{}born\textquotedblright{} (MW), and KRM gives the participle \emph{janita} for the root.
\end{cousinweb}

\begin{grammarlens}[title={What you've built}]
One more word for giving reasons and opinions.
\end{grammarlens}

\subsection*{Your turn}
\begin{itemize}
\raggedright
  \item \emph{Say it:} \emph{janayati}
  \item \emph{Say it:} \emph{janayati}, once more
  \item \emph{Say it:} \emph{krodhaḥ duḥkhaṃ janayati.} — Anger gives rise to sorrow.
  \item \emph{Recall:} say the Sanskrit for a lack, then the Sanskrit for comes about, then say \emph{janayati} again
\end{itemize}

\begin{tcolorbox}[breakable,colback=teal!4,colframe=teal!35!black,title={Before you move on}]
What does \sk{जनयति} mean? (\textquotedblleft{}Gives rise to; produces\textquotedblright{}.) Say it once more.
\end{tcolorbox}
